\chapter{Services métier et injection CDI}\label{chap:services}

\section{Pourquoi une couche de services}
Les servlets ne parlent jamais à JPA. Elles appellent quatre interfaces, qui décrivent \textit{ce que fait} l'application
sans dire \textit{comment}. Le tableau~\ref{tab:services} les présente.

\begin{table}[htbp]
\centering
\caption{Interfaces de service}\label{tab:services}
\begin{tabularx}{\linewidth}{@{}lX@{}}
\toprule
\textbf{Interface} & \textbf{Responsabilité}\\
\midrule
\code{InternauteService} & Inscrire un client, authentifier, retrouver un internaute.\\
\code{CatalogueService} & Lecture publique : vitrine active, catégories, produits par page de 9, recherche, fiche produit.\\
\code{GestionCatalogueService} & Administration : produits et catégories (création, modification, retrait, suppression).\\
\code{PanierService} & Consulter, ajouter, modifier une quantité, retirer, vider.\\
\bottomrule
\end{tabularx}
\end{table}

Ce découpage a deux effets. Les servlets restent courtes : lire les paramètres, appeler un service, choisir la
vue. Et les tests peuvent remplacer un service par une version simulée sans base de données (c'est ce que fait la classe
\code{ServiceChecks}).

\section{Transactions centralisées}\label{sec:transactions}
Sans serveur JTA, chaque opération doit ouvrir un \code{EntityManager}, une transaction, valider ou annuler, puis fermer.
Répété dans chaque méthode, ce code est une source d'oublis (une transaction jamais refermée, une connexion qui
fuit). Un composant unique, \code{TransactionRunner}, en est responsable. Il offre deux méthodes : \code{read} pour les
lectures, sans transaction, et \code{write} pour les écritures.

\lstinputlisting[language=Java,caption={TransactionRunner.write : validation, annulation, fermeture},basicstyle=\ttfamily\scriptsize]{extraits/TransactionRunner-write.txt}

Le \code{finally} garantit la fermeture quoi qu'il arrive. L'exception SQL brute n'est jamais montrée à
l'utilisateur : \code{SqlErrors.traduire} la convertit en message métier lisible. Les codes MySQL traités sont
1062 (doublon), 1451 et 1452 (clé étrangère), 1048 (valeur obligatoire manquante), 1406 (valeur trop longue), 3819 et
4025 (contrainte \code{CHECK}), ainsi que l'échec de verrouillage optimiste, reconnu par son nom de classe.

\section{Règles métier : à un seul endroit}
Les quantités, le stock et les produits inactifs sont contrôlés dans l'entité \code{Panier}, pas dans le service.
Le service se contente d'orchestrer : charger le client, charger le produit, appeler la méthode du panier, valider.

\lstinputlisting[language=Java,caption={Panier.ajouterProduit : fusion de lignes, prix figé},basicstyle=\ttfamily\scriptsize]{extraits/Panier-ajouterProduit.txt}

Un changement dans une règle (par exemple porter la quantité maximale de 99 à 50) se fait donc dans une constante unique.

\section{Validation des saisies}
La classe \code{Validation} collecte les erreurs au lieu de s'arrêter à la première. Le formulaire renvoyé à l'utilisateur
peut ainsi signaler tous les champs invalides d'un coup. Les contrôles couvrent le texte obligatoire, la longueur, le
format de l'e-mail, la complexité du mot de passe (8 caractères minimum, une lettre et un chiffre), le téléphone, la
référence produit (mise en majuscules), l'URL d'image (\code{https} uniquement) et les nombres décimaux ou entiers.
L'exception \code{ValidationException} porte la table des erreurs par champ.

\section{Requêtes}
Toutes les requêtes utilisent des paramètres liés, jamais de concaténation de saisies. Deux détails méritent d'être notés.

\begin{itemize}
  \item \textbf{Recherche par \code{LIKE}.} Les caractères \code{\%} et \code{\_} saisis par l'utilisateur sont
        échappés, avec \code{!} comme caractère d'échappement. Le choix de \code{!} et non de la barre oblique inverse est
        volontaire : sous MySQL, la barre oblique inverse a déjà un sens dans les chaînes littérales, ce qui rend
        l'échappement ambigu.
  \item \textbf{Requêtes dynamiques.} Les vérifications d'unicité sont construites selon les champs fournis, plutôt que
        d'utiliser \code{(:param IS NULL OR ...)}, une forme dont le comportement avec des paramètres nuls varie selon le pilote.
\end{itemize}

\section{Injection de dépendances avec CDI}
\subsection{Pourquoi CDI sur Tomcat}
Tomcat n'est pas un serveur d'applications complet : il ne fournit pas CDI. L'injection par \code{@Inject} est
pourtant la manière propre de relier servlets, services et gestionnaire de transactions sans \code{new} partout. L'implémentation
de référence, Weld, existe en version « servlet » qui s'embarque dans le WAR.

\subsection{Mise en place}
Trois fichiers suffisent.
\begin{itemize}
  \item \code{WEB-INF/beans.xml}, en mode \code{annotated} : seules les classes portant une annotation de portée
        deviennent des beans, ce qui limite l'analyse.
  \item \code{META-INF/context.xml} : déclare à Tomcat la ressource \code{BeanManager}.
  \item \code{WEB-INF/web.xml} : la référence \code{resource-env-ref} correspondante.
\end{itemize}
Au démarrage de Tomcat, la ligne \code{WELD-ENV-001100} (« Tomcat 7+ detected, CDI injection will be available in Servlets,
Filters and Listeners ») confirme que l'injection fonctionne.

\subsection{La fabrique d'EntityManager}
La fabrique d'\code{EntityManager} coûte cher à créer et peut être partagée entre les fils d'exécution ; l'\code{EntityManager}
lui-même ne le peut pas. Un producteur CDI crée donc \textbf{une} fabrique pour toute l'application, et un observateur
l'initialise dès le démarrage.

\lstinputlisting[language=Java,caption={EntityManagerFactoryProducer},basicstyle=\ttfamily\scriptsize]{extraits/EntityManagerFactoryProducer.txt}

L'initialisation dès le démarrage est une décision de confort : un mot de passe erroné ou une base arrêtée apparaît dans
la console de Tomcat à l'ouverture de l'application, et non à la requête du premier visiteur. La méthode annotée
\code{@Disposes} ferme la fabrique à l'arrêt, ce qui évite de laisser des connexions ouvertes lors d'un redéploiement.
